%% hopfion-framework/quantum/bell.tex
%% Sources: main_paper8.tex, main_paper9.tex
%% Geometry-dependent Bell model (VIII), Born rule, Malus law, three CHSH regimes, Hurwitz algebras (IX)
%% Auto-generated from project files. Edit source papers to update.

%════════════════════════════════════════════════════════════════════
% Geometry-Dependent Bell Correlations — Paper VIII
%════════════════════════════════════════════════════════════════════

% Four axioms of the density-feedback Hopfion framework, recalled from
% Paper V (derived in Papers I--VI).
\begin{equation}
  S_{\mathrm{eff}}(\theta, \rho)
  \;=\; \frac{\sin^4\!\theta}{\vp^6\,(1 + \bst\,\rho)},
  \label{P8:eq:seff_axiom}
\end{equation}
\status{published}
\source{P8:eq:seff_axiom}

\begin{equation}
  M_n = M_0\,\vp^{2n}, \qquad
  m_f = m_0\,\vp^{-2n_f},
  \label{P8:eq:hierarchy_axiom}
\end{equation}
\status{published}
\source{P8:eq:hierarchy_axiom}

\begin{equation}
  S_{\mathrm{parent}}[\rho]
  \;=\; \int\!\mathrm{d}^4x\,\sqrt{-g}\;
  \frac{|\nabla\rho|^2}{\vp^6\bigl(1 + \bst\,|\nabla\rho|^2\bigr)},
  \label{P8:eq:parent_action}
\end{equation}
\status{published}
\source{P8:eq:parent_action}

\begin{remark}[Extension of the density variable]
\label{P8:rem:rho_extension}
In Papers~I--IV the feedback denominator $1+\bst\kappa$ uses the local
field curvature $\kappa$ (the FN gradient density) as the density variable.
In this paper, the parent action of Paper~I~\cite{Paper1} uses a scalar
field $\rho$ as an independent dynamical variable, with $1+\beta\rho$ in
the denominator.  The two descriptions are compatible: in the static
condensate background, $\rho\propto\kappa$ and $\beta\propto\bst$, so
Axiom~A2 recovers the Papers~I--IV form.  The $\rho$-based description
is a $k$-essence extension is used for the analysis in this paper.
\end{remark}
\status{published}
\source{P8:rem:rho_extension}

%──────────────────────────────────────────────────────────────────────────────
\subsection*{The Condensate Orientation as Correlation Variable}

Consider two spatially separated measurement sites A and B (Alice and Bob). The pair creation event fixes the local condensate gradient orientation $\hat{n} = \nabla\rho/|\nabla\rho|$ at the source, which determines the correlation structure through the suppression law. This direction is randomly distributed over the sphere but fixed for each pair.

Alice measures along direction $\hat{a}$ (angle $\alpha$) and Bob along $\hat{b}$ (angle $\beta$). The mismatch angles between the measurement axes and the condensate orientation $\hat{n}$ are
\begin{equation}
\theta_A = \arccos(\hat{a} \cdot \hat{n}), \quad \theta_B = \arccos(\hat{b} \cdot \hat{n}).
\label{P8:eq:mismatch_angles}
\end{equation}

The orientation $\hat{n}$ is set by the condensate density profile at the source and determines the correlation structure through the suppression law $S(\theta)=\sin^4\!\theta/\vp^6$.
\status{published}
\source{P8:eq:mismatch_angles}

%──────────────────────────────────────────────────────────────────────────────
\subsection*{Measurement-Induced Perturbations}

Crucially, measurement is not a passive observation. When Alice's detector interacts with the system, it creates a localized density perturbation $\delta\rho_A > 0$ (from energy deposition, condensate back-action, or quantum interaction). This perturbation modifies the effective suppression at A:
\begin{equation}
S_{\mathrm{eff}}(\theta_A, \rho + \delta\rho_A) = \frac{1}{\phi^6} \sin^4\theta_A \Big/ (1 + \beta (\rho + \delta\rho_A)).
\label{P8:eq:seff_perturbed}
\end{equation}

The denominator $(1 + \beta(\rho + \delta\rho_A))$ becomes larger when $\delta\rho_A > 0$, making $S_{\mathrm{eff}}$ smaller (weaker suppression). This creates a preferential low-suppression channel for condensate orientations aligned with the measurement axis $\hat{a}$. Similarly for Bob at site B.

Importantly, $\delta\rho_A$ and $\delta\rho_B$ are direction-dependent: stronger when $\hat{n}$ aligns with the measurement axis (small $\theta$), weaker for misaligned directions. This selective softening enhances correlations for relevant hidden variables without requiring non-locality---the measurement locally reshapes the condensate density field.
\status{published}
\source{P8:eq:seff_perturbed}

%──────────────────────────────────────────────────────────────────────────────
\subsection*{Probability Weights from Suppression}

The probability of outcome "+1" at Alice's detector (rebalancing along $\hat{a}$) is determined by the suppression cost. The unnormalized weight is
\begin{equation}
\tilde{w}_A(+1 \mid \hat{n}) = \exp\left( -\lambda S_{\mathrm{eff}}(\theta_A, \rho + \delta\rho_A) \right),
\label{P8:eq:weight_plus}
\end{equation}
where $\lambda > 0$ is a coupling strength. For the opposite outcome "-1" (effective angle $\theta' = \pi/2 - \theta_A$, so $\sin\theta' = \cos\theta_A$):
\begin{equation}
\tilde{w}_A(-1 \mid \hat{n}) = \exp\left( -\lambda \cdot \frac{1}{\phi^6} \cos^4\theta_A \Big/ (1 + \beta (\rho + \delta\rho_A)) \right).
\label{P8:eq:weight_minus}
\end{equation}

Normalizing over outcomes gives the probability
\begin{equation}
p_A(+1 \mid \hat{n}) = \frac{\tilde{w}_A(+1)}{\tilde{w}_A(+1) + \tilde{w}_A(-1)}.
\label{P8:eq:prob_plus}
\end{equation}

For small angles ($\theta_A \ll 1$), $\sin^4\theta_A \approx \theta_A^4$, and the normalized probability approximates
\begin{equation}
p_A(+1 \mid \theta_A) \approx \left[ \cos^4\left(\frac{\theta_A}{2}\right) \right]^{f_A},
\label{P8:eq:prob_smallangle}
\end{equation}
where $f_A = f_{\mathrm{base}} / (1 + \beta \cos\theta_A)$ accounts for direction-dependent density feedback (lower $f_A$ when $\theta_A \approx 0$ due to measurement-induced softening). The exponent 4 (from $\sin^4\theta$) produces the characteristic quantum form for spin-1/2 fermions. Similarly for Bob with $p_B(+1 \mid \hat{n})$ and softening parameter $f_B$.
\status{published}
\source{P8:eq:weight_plus,P8:eq:weight_minus,P8:eq:prob_plus,P8:eq:prob_smallangle}

%──────────────────────────────────────────────────────────────────────────────
\subsection*{Expectation Values and CHSH}

For singlet-like anti-correlated pairs (condensate orientations $\hat{n}$ and $-\hat{n}$ at each site), the correlation expectation value for measurement angles $\alpha$ and $\beta$ is obtained by averaging over all possible hidden-variable directions $\hat{n}$:
\begin{equation}
E(\alpha, \beta) = \frac{1}{4\pi} \int d\Omega \, [p_{\mathrm{same}}(\hat{n}) - p_{\mathrm{diff}}(\hat{n})],
\label{P8:eq:E_expectation}
\end{equation}
where
\begin{align}
p_{\mathrm{same}} &= p_A(+1)(1 - p_B(+1)) + (1 - p_A(+1))p_B(+1), \\
p_{\mathrm{diff}} &= p_A(+1)p_B(+1) + (1 - p_A(+1))(1 - p_B(+1)).
\label{P8:eq:p_same_diff}
\end{align}

The CHSH correlator for standard Bell test angles ($a=0^\circ$, $a'=45^\circ$, $b=22.5^\circ$, $b'=67.5^\circ$) is
\begin{equation}
\mathrm{CHSH} = |E(a,b) + E(a,b')| + |E(a',b) - E(a',b')|.
\label{P8:eq:chsh_def}
\end{equation}

Classical local hidden-variable theories obey $\mathrm{CHSH} \leq 2$. Quantum mechanics predicts $\mathrm{CHSH} = 2\sqrt{2} \approx 2.828$ (Tsirelson's bound~\cite{Tsirelson1980}) for maximally entangled states.
\status{published}
\source{P8:eq:E_expectation,P8:eq:p_same_diff,P8:eq:chsh_def}

%──────────────────────────────────────────────────────────────────────────────
\subsection*{Dynamic Rescaling Factor}

To account for the geometric dilution effect quantitatively, we introduce a dynamic rescaling factor based on the actual sampled geometry. The raw correlation $E_{\mathrm{raw}}$ computed from Eq. (14) is diluted by the effective solid-angle weight $\langle \sin^4\theta \rangle$ over the sampled $\hat{n}$ distribution. The rescaled correlation is
\begin{equation}
E_{\mathrm{rescaled}} = \frac{E_{\mathrm{raw}}}{\langle \sin^4\theta \rangle_{\mathrm{geometry}}},
\label{P8:eq:E_rescaled}
\end{equation}
where
\begin{equation}
\langle \sin^4\theta \rangle_{\mathrm{geometry}} = \frac{1}{N} \sum_{i=1}^N \sin^4\theta_i
\label{P8:eq:sin4_geometry_avg}
\end{equation}
is computed from the Monte Carlo sampled hidden-variable directions.

For pure 2D plane-constrained geometry, $\langle \sin^4\theta \rangle = 3/8 = 0.375$ exactly
(proved analytically in Paper~IX~\cite{Paper9}, Step~1 of the Tsirelson gap theorem:
$(1/\pi)\int_0^\pi\sin^4\theta\,d\theta = 3/8$), giving a rescale factor of $8/3\approx2.667$.
For full 3D isotropic sampling,
$\langle \sin^4\theta \rangle = 8/15 \approx 0.533$ exactly
(spherical average $(1/2)\int_0^\pi\sin^4\theta\sin\theta\,d\theta = 8/15$),
giving a rescale factor of $15/8 = 1.875$.
For biased 3D geometries (Gaussian distribution around mid-plane), intermediate values arise.
This rescaling is not a fit parameter but a geometric correction determined by the sampled
distribution---it accounts for the dilution inherent in averaging over misaligned directions.
\status{published}
\source{P8:eq:E_rescaled,P8:eq:sin4_geometry_avg}
\depends{quantum/bell.tex (P9:thm:tsirelson_gap)}

%──────────────────────────────────────────────────────────────────────────────
\subsection*{Plane-Constrained Configuration}

In the 2D plane-constrained case, the condensate orientation $\hat{n}$ is restricted to the xy-plane (perpendicular to the z-axis). Alice and Bob's measurement axes $\hat{a}$ and $\hat{b}$ also lie in this plane. For standard CHSH angles ($a=0^\circ$, $a'=45^\circ$, $b=22.5^\circ$, $b'=67.5^\circ$), all angles are coplanar.

The solid-angle integration in Eq. (14) reduces to an azimuthal integral over $\phi \in [0, 2\pi]$:
\begin{equation}
E(\alpha, \beta) = \frac{1}{2\pi} \int_0^{2\pi} [p_{\mathrm{same}}(\phi) - p_{\mathrm{diff}}(\phi)] \, d\phi,
\label{P8:eq:E_planar}
\end{equation}
where $\hat{n} = (\cos\phi, \sin\phi, 0)$ samples directions uniformly in the plane.
\status{published}
\source{P8:eq:E_planar}

%──────────────────────────────────────────────────────────────────────────────
\subsection*{Results: Approaching Tsirelson's Bound}

Table 1 shows CHSH values for plane-constrained geometry across different density feedback strengths $\beta$. For mild to moderate feedback ($\beta = 0.1$--$1.0$), the rescaled CHSH ranges from 2.57 to 2.62, exceeding the classical bound of 2 and approaching Tsirelson's quantum bound $2\sqrt{2} \approx 2.828$ within $\sim 7$--$9\%$. The Hopfion condensate framework fixes $\beta^* \approx 0.452$ from first principles (Paper~I~\cite{Paper1}); this row is highlighted in the table.

\begin{table}[h]
\centering
\caption{CHSH values in plane-constrained geometry with dynamic rescaling.
Density feedback coupling $\beta$ ranges from minimal (0.1) to moderate (1.0),
computed using the same direction-dependent softening method throughout.
The highlighted row $\beta^* = 0.452$ is the value fixed by the Hopfion
condensate framework (Paper~I~\cite{Paper1}); its CHSH falls between $\beta=0.3$ and
$\beta=0.5$ as expected. Monte Carlo, N=100,000 samples, $f_{\mathrm{base}}=2.0$.}
\label{P8:tab:planar_chsh}
\begin{tabular}{c|c|c}
\hline
$\beta$ & Raw CHSH & Rescaled CHSH \\
\hline
0.1 & 0.982 & 2.624 \\
0.3 & 0.980 & 2.622 \\
$\boldsymbol{\beta^* \approx 0.452}$ & $\mathbf{0.978}$ & $\mathbf{2.611}$ \\
0.5 & 0.977 & 2.596 \\
0.7 & 0.977 & 2.619 \\
1.0 & 0.969 & 2.573 \\
\hline
\end{tabular}
\end{table}

The results exhibit remarkable stability across the physically relevant range $\beta = 0.1$--$1.0$, varying by only $\sim 2\%$ ($\Delta\mathrm{CHSH} \approx 0.05$). This robustness indicates the violation is not a fine-tuned artifact but emerges naturally from the angular suppression structure combined with measurement-induced softening in planar geometries.

The geometric rescale factor is dynamically computed from the actual sampled geometry and remains approximately constant across $\beta$ values (varying between 2.57 and 2.62 for planar configuration), as expected---it depends only on $\langle \sin^4\theta \rangle$ over the sampled distribution, not on the density feedback strength. This confirms the rescaling is a genuine geometric correction, not a hidden fit parameter.
\status{published}
\source{P8:tab:planar_chsh}

%──────────────────────────────────────────────────────────────────────────────
\subsection*{The Framework-Fixed Coupling $\beta^* \approx 0.452$}

The Hopfion condensate framework of Papers~I--VII fixes the density-feedback
coupling uniquely from the golden-ratio self-consistency condition
$J_{2\mathrm{iso}}^{\mathrm{fb}}/J_{2a} = \varphi$ at the condensate saddle (Paper~I~\cite{Paper1}):
\begin{equation}
  \beta^* \;\approx\; 0.452.
  \label{P8:eq:beta_star}
\end{equation}
This is not a free parameter of the Bell model.
It is the same value that determines the lepton mass hierarchy, the Bogomolny
parameter $\lambda = \varphi^6$, and the condensate sound speed $c_s = c/\varphi$
in Papers~I--VII.

The highlighted row in Table~1 shows that at $\beta^*$, using the same
direction-dependent softening method as all other rows, the model gives
CHSH$_\mathrm{raw} = 0.978$, CHSH$_\mathrm{rescaled} = 2.611$ --- correctly
interpolating between the $\beta=0.3$ and $\beta=0.5$ rows.
The companion Paper~IX~\cite{Paper9} also analyses a simpler limiting case (pure
$f=2$, no direction dependence, no density softening) which gives
CHSH$_\mathrm{rescaled}=2.361$ at $\beta^*$; the gap between that value and
the full-model $2.611$ is the measurement-strength enhancement from the
direction-dependent $f$ used here.
\status{published}
\source{P8:sec:beta_star}
\depends{quantum/bell.tex (P8:tab:planar_chsh)}

%──────────────────────────────────────────────────────────────────────────────
\subsection*{Biased 3D Geometries: Smooth Transition}

To understand the transition from planar to fully isotropic 3D geometry, we also evaluate biased configurations where $\hat{n}$ is sampled from a Gaussian distribution centered on the xy-plane with width $\sigma$. Smaller $\sigma$ corresponds to stronger planar bias; larger $\sigma$ approaches isotropic sampling.

\begin{table}[h]
\centering
\caption{CHSH values for biased 3D geometries at the framework-fixed coupling
$\beta^*\approx0.452$. Bias parameter $\sigma$ controls out-of-plane spread:
$\sigma=0$ is pure 2D plane, $\sigma \to \infty$ is full 3D isotropic.
Results at $\beta=0.5$ differ by $<0.002$, confirming the robustness
of the geometry dependence.}
\label{P8:tab:biased3d_chsh}
\begin{tabular}{c|c|c}
\hline
Geometry & Raw CHSH & Rescaled CHSH \\
\hline
2D plane ($\sigma = 0$) & 0.981 & 2.613 \\
Biased 3D ($\sigma = \pi/18 \approx 10^\circ$) & 0.975 & 2.535 \\
Biased 3D ($\sigma = \pi/12 \approx 15^\circ$) & 0.972 & 2.454 \\
Biased 3D ($\sigma = \pi/6 \approx 30^\circ$) & 0.943 & 2.022 \\
3D isotropic ($\sigma \to \infty$) & 0.903 & 1.692 \\
\hline
\end{tabular}
\end{table}

The transition is smooth and monotonic: as geometry becomes more isotropic (increasing $\sigma$), CHSH degrades from quantum-like ($\approx 2.6$) toward classical ($\approx 1.7$). Even modest out-of-plane spread ($\sigma = \pi/6 \approx 30^\circ$) reduces CHSH to 2.02, barely exceeding the classical bound. This demonstrates that 3D dilution is not a binary effect but a continuous function of geometric bias, providing multiple intermediate points for experimental validation.

The near-identical results at $\beta^* = 0.452$ and $\beta = 0.5$ (differences $< 0.002$ across all geometries) confirm that the geometry-dependence curve is not a property of any particular $\beta$ value but of the angular suppression structure $\sin^4\theta/\vp^6$ itself. The algebraic reason for this robustness is identified in Paper~IX~\cite{Paper9}: the $\sin^4\theta$ weight is the quaternionic ($\mathbb{H}$-valued) norm on transition amplitudes, and the Fourier harmonic structure of the resulting correlation function is $\beta$-independent at leading order.
\status{published}
\source{P8:tab:biased3d_chsh}

%──────────────────────────────────────────────────────────────────────────────
\subsection*{Full Sphere Sampling}

In the 3D isotropic case, $\hat{n}$ is sampled uniformly over the sphere using the standard method: $\cos\theta \sim \mathrm{Uniform}(-1, 1)$ and $\phi \sim \mathrm{Uniform}(0, 2\pi)$, giving $\hat{n} = (\sin\theta\cos\phi, \sin\theta\sin\phi, \cos\theta)$. Alice and Bob's measurement axes $\hat{a}$ and $\hat{b}$ remain in the xy-plane for simplicity, but the hidden variable now explores all directions including those perpendicular to the measurement plane.

The solid-angle integration in Eq. (14) becomes a full spherical average:
\begin{equation}
E(\alpha, \beta) = \frac{1}{4\pi} \int_0^{2\pi} d\phi \int_0^\pi \sin\theta \, d\theta \, [p_{\mathrm{same}}(\theta, \phi) - p_{\mathrm{diff}}(\theta, \phi)].
\label{P8:eq:E_isotropic}
\end{equation}

For $\hat{n}$ oriented out of the measurement plane (large $\theta$ relative to z-axis), the mismatch angles $\theta_A$ and $\theta_B$ to measurement axes $\hat{a}$ and $\hat{b}$ become nearly random, washing out the $\sin^4\theta$ angular structure. This dilution effect is the core prediction distinguishing this model from quantum mechanics.
\status{published}
\source{P8:eq:E_isotropic}

%──────────────────────────────────────────────────────────────────────────────
\subsection*{Non-Perturbative Form: The Unobservable Baseline}

To understand the pre-measurement state, we evaluate the model without measurement-induced perturbations: $\delta\rho_A = \delta\rho_B = 0$. In this non-perturbative regime, only the background condensate density $\rho$ provides feedback via the denominator $(1 + \beta\rho)$. The probability weights become
\begin{equation}
\tilde{w}_A(\pm 1 \mid \hat{n}) = \exp\left( -\lambda \cdot \frac{1}{\phi^6} \frac{\sin^4\theta_A \text{ or } \cos^4\theta_A}{1 + \beta\rho} \right),
\label{P8:eq:weight_nonperturbative}
\end{equation}
with no direction-dependent softening (no $\delta\rho$ to create preferential channels).
\status{published}
\source{P8:eq:weight_nonperturbative}

%──────────────────────────────────────────────────────────────────────────────
\subsection*{Results: Classical Regime}

Table 3 shows CHSH values for fully isotropic 3D geometry in the non-perturbative regime across $\beta = 0.1$--$2.0$. The rescaled CHSH ranges from 1.03 (minimal baseline softening) down to 0.39 (strong softening), all significantly below the classical bound of 2.0.

\begin{table}[h]
\centering
\caption{CHSH values for isotropic 3D geometry in non-perturbative regime ($\delta\rho = 0$). Baseline density feedback $\beta$ ranges from minimal (0.1) to strong (2.0). Monte Carlo evaluation with N=100,000 samples, dynamic rescaling.}
\label{P8:tab:isotropic_nonpert_chsh}
\begin{tabular}{c|c|c|c}
\hline
$\beta$ & Raw CHSH & Rescaled CHSH & Rescale Factor \\
\hline
0.1 & 0.552 & 1.034 & 1.874 \\
0.3 & 0.507 & 0.949 & 1.873 \\
0.5 & 0.465 & 0.873 & 1.877 \\
0.7 & 0.419 & 0.786 & 1.877 \\
1.0 & 0.354 & 0.662 & 1.871 \\
2.0 & 0.210 & 0.395 & 1.876 \\
\hline
\end{tabular}
\end{table}

These values represent the fundamental pre-measurement state of the system. The smooth degradation with increasing $\beta$ arises because stronger baseline softening washes out the angular structure---when suppression is uniformly weak everywhere, there's no preferential direction, leading to nearly random correlations.

The rescale factor remains nearly constant ($\sim 1.87$) across all $\beta$ values, as expected for geometry-dependent correction. It differs slightly from the planar value due to the full 3D solid-angle average having $\langle \sin^4\theta \rangle \approx 0.6$--$0.65$ (compared to 0.533 for pure 2D), but the consistency confirms the rescaling is a genuine geometric effect, not a tuning parameter.
\status{published}
\source{P8:tab:isotropic_nonpert_chsh}

%──────────────────────────────────────────────────────────────────────────────
\subsection*{Comparison: 3D Isotropic with Measurement}

For completeness, Table 4 shows CHSH values in 3D isotropic geometry with measurement-induced perturbations (the code from Section 4 applied to full sphere sampling). Even with $\delta\rho > 0$, the 3D dilution effect remains dominant, yielding CHSH $\approx$ 1.7 (barely classical).

\begin{table}[h]
\centering
\caption{CHSH values for isotropic 3D geometry with measurement perturbations ($\delta\rho > 0$, $f_{\mathrm{base}}=2.0$). Compared to Table 3 (non-perturbative), measurement softening boosts CHSH but cannot overcome 3D dilution.}
\label{P8:tab:isotropic_measured_chsh}
\begin{tabular}{c|c|c}
\hline
$\beta$ & Raw CHSH & Rescaled CHSH \\
\hline
0.1 & 0.911 & 1.711 \\
0.3 & 0.902 & 1.693 \\
0.5 & 0.900 & 1.687 \\
0.7 & 0.900 & 1.691 \\
1.0 & 0.894 & 1.682 \\
\hline
\end{tabular}
\end{table}

Measurement increases CHSH from ~1.0 (non-perturbative, Table 3) to ~1.7 (with perturbation, Table 4), demonstrating that $\delta\rho$ does enhance correlations. However, the boost is insufficient to reach quantum-like violations (~2.6) because out-of-plane $\hat{n}$ directions still contribute near-random outcomes, diluting the angular structure. Only in planar or biased geometries, where $\hat{n}$ is constrained to relevant directions, does measurement-induced softening produce strong violations.

This is the key testable prediction: this model predicts CHSH $\approx$ 1.7 in true 3D isotropic experiments (even with standard detectors creating $\delta\rho > 0$), while quantum mechanics predicts CHSH = 2.828 independent of geometry.
\status{published}
\source{P8:tab:isotropic_measured_chsh}
\residual{P (untested; falsifiable via 3D-isotropic Bell test, Section~5)}

%──────────────────────────────────────────────────────────────────────────────
\subsection*{Summary: Geometry-Dependent Predictions}

The Density-Feedback Hopfion model makes three clear, testable predictions regarding geometry:

\begin{enumerate}
\item \textbf{Planar geometry} (all existing experiments): CHSH $\approx$ 2.6, approaching Tsirelson's bound. Consistent with observations.

\item \textbf{Biased 3D geometry} (slight out-of-plane spread): CHSH degrades smoothly from 2.6 (pure planar) to 2.0 ($\sigma \approx 30^\circ$) as shown in Table 2. Testable with controlled geometric bias.

\item \textbf{Isotropic 3D geometry} (full sphere, no bias): CHSH $\approx$ 1.7 with measurement (Table 4), or CHSH $\approx$ 1.0 for weak measurements (approaching Table 3 non-perturbative values). Significantly below classical bound in weak-measurement limit.
\end{enumerate}

Standard quantum mechanics predicts CHSH = 2.828 for all three cases (up to projection factors). This geometry dependence provides a clear experimental distinction, testable with next-generation satellite networks or 3D entangled-state protocols (Section 6).
\status{P}
\source{Paper VIII, sec. Summary: Geometry-Dependent Predictions}
\depends{quantum/bell.tex (P8:tab:planar_chsh, P8:tab:biased3d_chsh, P8:tab:isotropic_nonpert_chsh, P8:tab:isotropic_measured_chsh)}

%════════════════════════════════════════════════════════════════════
% Born Rule and Observable Structures — Paper IX
%════════════════════════════════════════════════════════════════════
\begin{theorem}[Hopf push-forward]
\label{P9:thm:hopf_pushforward}
The condensate signed correlation $\xcond$ is the push-forward of the
standard QM signed correlation $x_\mathrm{QM}$ under the Hopf fibration:
\begin{equation}
  \xcond(\theta) \;=\; x_\mathrm{QM}(\eta^{-1}_\mathrm{avg}(\theta)),
  \label{P9:eq:hopf_pushforward}
\end{equation}
or equivalently,
\begin{equation}
  \xcond(\theta) \;=\; \frac{2\cos\theta}{1+\cos^2\theta}
  \;=\; \frac{x_\mathrm{QM}}{1 - \tfrac{1}{2}(1-x_\mathrm{QM}^2)},
  \label{P9:eq:xcond_from_xQM}
\end{equation}
where the denominator $1+\cos^2\theta = 2-\sin^2\theta$ is the
$S^3$-to-$S^2$ Jacobian factor in the Hopf fibration.
\end{theorem}
\status{published}
\source{P9:thm:hopf_pushforward}

\begin{theorem}[Exact Fourier series of $\xcond$]
\label{P9:thm:fourier_exact}
\begin{equation}
  \xcond(\theta)
  \;=\; \sum_{k=0}^{\infty} A_{2k+1}\,\cos((2k+1)\theta),
  \label{P9:eq:xcond_fourier}
\end{equation}
with \emph{exact} coefficients
\begin{equation}
  \boxed{A_{2k+1}
  \;=\; (4-2\sqrt{2})\cdot(-(3-2\sqrt{2}))^k
  \;=\; (4-2\sqrt{2})\cdot(2\sqrt{2}-3)^k},
  \label{P9:eq:A_exact}
\end{equation}
and the Parseval identity
\begin{equation}
  \boxed{\sum_{k=0}^{\infty} \frac{A_{2k+1}^2}{2}
  \;=\; C(0) \;=\; \frac{1}{\sqrt{2}}}.
  \label{P9:eq:parseval_exact}
\end{equation}
In particular $A_1 = 4-2\sqrt{2} = 2(2-\sqrt{2})\approx1.172$.
\end{theorem}
\status{published}
\source{P9:thm:fourier_exact}

\begin{theorem}[Tsirelson gap decomposition]
\label{P9:thm:tsirelson_gap}
The CHSH of the condensate model in the 2D planar geometry, after applying the
geometric rescaling $\CHSH_\mathrm{rescaled}=\CHSH_\mathrm{raw}/\langle\sin^4\theta\rangle_\mathrm{2D}$
with $\langle\sin^4\theta\rangle_\mathrm{2D} = 3/8$, satisfies
\begin{equation}
  \boxed{\CHSH_\mathrm{rescaled}
  \;=\; \frac{8}{3}\,\CHSH_\mathrm{raw}
  \;\approx\; 2.36},
  \label{P9:eq:CHSH_rescaled_value}
\end{equation}
and the gap from Tsirelson's bound decomposes as
\begin{equation}
  \Tsirelson - \CHSH_\mathrm{rescaled}
  \;\approx\; 0.47
  \;=\; \underbrace{\vphantom{X}\Tsirelson(1-A_1)}_{\text{harmonic mismatch}}
  \;+\; \underbrace{\vphantom{X}\Tsirelson\!\left(A_1 - \frac{3}{8}\cdot\frac{8}{3}\right)}_{\text{projection error}},
  \label{P9:eq:gap_decomposition}
\end{equation}
where $A_1 = 4-2\sqrt{2}\approx1.172$ is the leading Fourier coefficient
of $\xcond$ (Theorem~\ref{P9:thm:fourier_exact}, equation~\eqref{P9:eq:A_exact}).
 
Under the pointwise $\HH\to\CC$ projection defined by replacing
$\pcond(\theta)\to\pQM(\theta)$ (i.e., $t^4\to t^2$ in the denominator),
\begin{equation}
  \CHSH_\mathrm{projected} \;=\; \Tsirelson.
  \label{P9:eq:CHSH_projected_exact}
\end{equation}
\end{theorem}
\status{published}
\source{P9:thm:tsirelson_gap}

\begin{theorem}[Born rule from complexification]
\label{P9:thm:born_rule}
Define the $\CC$-valued amplitude by complexification via the topological
winding phase:
\begin{equation}
  \psi_\CC(\theta)
  \;=\; |\psi_\HH(\theta)|_\HH^{1/2} \cdot e^{i\Phi(\theta)},
  \label{P9:eq:complexification}
\end{equation}
where $\Phi(\theta) = Q\cdot\omega(\theta)$, $Q=10$ is the Hopf charge,
and $\omega(\theta) = 2\pi(1-\cos\theta)/2$ is the solid angle subtended
at polar angle $\theta$.
Then the Born rule holds exactly:
\begin{equation}
  \boxed{P(\theta) \;=\; |\psi_\CC(\theta)|^2
  \;=\; |\psi_\HH(\theta)|_\HH
  \;=\; \frac{e^{-\sin^4\theta/(2\vp^6)}}{Z(\theta)^{1/2}}}.
  \label{P9:eq:born_from_condensate}
\end{equation}
In the strong-measurement limit $(\delta\rho/\rho\to\infty$,
Section~3 of~\cite{Paper8}$)$, this converges to the
standard Born rule:
\begin{equation}
  P(\theta)
  \;\xrightarrow{\delta\rho\to\infty}\;
  \frac{\cos^4(\theta/2)}{\cos^4(\theta/2)+\sin^4(\theta/2)}
  \;=\; \pcond(\theta)
  \;\approx\; \pQM(\theta)\bigl[1 + O(\sin^4\theta/\vp^6)\bigr].
  \label{P9:eq:born_limit}
\end{equation}
\end{theorem}
\status{published}
\source{P9:thm:born_rule}

\begin{remark}[Division algebra assignments]
\label{P9:rem:hurwitz_roles}
The three non-real Hurwitz algebras each play a distinct role in the
derivation of the Born rule from the condensate:
\begin{enumerate}[label=(\roman*)]
  \item $\HH$ \emph{(quaternions, second generation lepton $k=2$):}
    supplies the $|\cdot|^4$ suppression magnitude
    $\sin^4\theta = |\langle\hat{a}\mid\hat{n}\rangle_\perp|_\HH^4$,
    driving the Bell correlations of~\cite{Paper8}.

  \item $\CC$ \emph{(complex numbers, first generation lepton $k=1$):}
    supplies the winding phase $e^{i\Phi}$ in~\eqref{P9:eq:complexification},
    enabling interference, the inner product structure, and the tensor
    product structure for composite entangled systems.
    Standard quantum mechanics is $J_n(\CC)$ because only over $\CC$
    does $(A\otimes B)_\CC$ remain a valid $\CC$-Hilbert space.

  \item $\OO$ \emph{(octonions, third generation lepton $k=3$):}
    supplies the exceptional Albert algebra $J_3(\OO)$ whose automorphism
    group $G_2\subset F_4$ constrains the Hopf charge to
    $Q\in\ZZ = \pi_3(S^2)$ (not $\pi_3(S^3)=\ZZ$ for generic spheres).
    The non-associativity of $\OO$ forces the three-generation structure
    and prevents the Cayley--Dickson tower from continuing beyond $\OO$.
\end{enumerate}
The Born rule is the interface between (i) and (ii): $\HH$ provides the
magnitude; $\CC$ provides the phase; $P=|\psi_\CC|^2$ is the result.
\end{remark}
\status{published}
\source{P9:rem:hurwitz_roles}

\begin{theorem}[$n$-particle GHZ Mermin correlator]
\label{P9:thm:GHZ}
Let $\hat{n}\in S^2$ be the shared hidden-variable direction, and let
$\theta_i = \arccos(\hat{a}_i\cdot\hat{n})$ be the mismatch angle for
particle $i$, $i=1,\ldots,n$.
The $n$-particle condensate Mermin correlator in the 2D planar geometry is
\begin{equation}
  M_\mathrm{raw}^{(n)}(\hat{a}_1,\ldots,\hat{a}_n)
  \;=\; \left\langle\prod_{i=1}^n \xcond(\theta_i)\right\rangle_{\!\hat{n}},
  \label{P9:eq:Mermin_raw}
\end{equation}
where $\xcond(\theta) = 2\cos\theta/(1+\cos^2\theta)$.
Under the pointwise $\HH\to\CC$ projection ($\xcond\to x_\mathrm{QM}=\cos\theta$),
the $n$-particle Mermin value reaches the quantum Tsirelson--Mermin
bound exactly:
\begin{equation}
  M_\mathrm{projected}^{(n)}
  \;=\; \left\langle\prod_{i=1}^n \cos\theta_i\right\rangle_{\!\hat{n},\,\mathrm{opt}}
  \;=\; 2^{n/2},
  \label{P9:eq:Mermin_QM}
\end{equation}
where ``opt'' denotes the optimal GHZ measurement angles.
The $n$-particle autocorrelation satisfies the \emph{$n$-fold Parseval identity}:
\begin{equation}
  \boxed{C^{(n)}(0)
  \;\equiv\; \left\langle\xcond(\theta)^{2n}\right\rangle_{\!\theta=0}
  \;=\; \left(\frac{1}{\sqrt{2}}\right)^n}.
  \label{P9:eq:parseval_n}
\end{equation}
\end{theorem}
\status{published}
\source{P9:thm:GHZ}

\begin{theorem}[CHSH$(f)$ transition]
\label{P9:thm:chsh_f}
Let $\CHSH_\mathrm{raw}(f)$ be the raw CHSH value of the $f$-parameterised
model~\eqref{P9:eq:p_f} in the 2D planar geometry at the standard optimal angles,
and let $\CHSH_f = \CHSH_\mathrm{raw}(f)/\langle\sin^{2f}\theta\rangle_\mathrm{2D}$
be the $f$-rescaled CHSH.
 
\begin{enumerate}[label=\emph{(\roman*)}]
  \item \emph{(Monotonicity.)}
    $\CHSH_\mathrm{raw}(f)$ is strictly increasing in $f$ for $f\geq1$:
    \begin{equation}
      \CHSH_\mathrm{raw}(1) < \CHSH_\mathrm{raw}(2) < \cdots
      \;\to\; 1 \quad\text{as }f\to\infty.
      \label{P9:eq:CHSH_monotone}
    \end{equation}
 
  \item \emph{(Exact values.)}
    \begin{align}
      \CHSH_\mathrm{raw}(f)
      &= \bigl|C_f(\pi/8) + C_f(3\pi/8)\bigr|,
      \label{P9:eq:CHSH_exact_f}\\
      C_f(\psi)
      &= \frac{1}{2\pi}\int_0^{2\pi} x_f(\varphi)\,x_f(\varphi+\psi)\,d\varphi,
      \nonumber
    \end{align}
    with $\CHSH_\mathrm{raw}(1) = \tfrac{1}{2}(\cos(\pi/8)+\cos(3\pi/8))
    = \tfrac{1}{2}\sqrt{1+\tfrac{1}{\sqrt{2}}}\,$.
 
  \item \emph{(Tsirelson recovery.)}
    Under the $\HH\to\CC$ projection ($f=2\to f=1$) with $f$-rescaling:
    \begin{equation}
      \CHSH_1
      = \frac{\CHSH_\mathrm{raw}(1)}{1/2}
      = \sqrt{1+\tfrac{1}{\sqrt{2}}} \;\approx\; 1.307,
      \label{P9:eq:CHSH_f1_rescaled}
    \end{equation}
    and under the normalisation identification that maps
    $E_\mathrm{raw}(\psi) = -C_1(\psi)$ to the standard QM convention
    $E_\mathrm{QM}(\psi) = -\cos\psi$ (a factor-$2$ renormalisation of
    the correlation function):
    \begin{equation}
      \CHSH_\mathrm{QM} \;=\; 2\cdot\CHSH_1
      \;=\; 2\sqrt{1+\tfrac{1}{\sqrt{2}}} \;\approx\; 2.613.
      \label{P9:eq:norm_id}
    \end{equation}
    The Tsirelson bound $\Tsirelson\approx2.828$ is recovered precisely
    when the $\HH\to\CC$ projection is combined with the standard-angle
    optimisation over all $(\alpha,\alpha',\beta,\beta')$~\cite{Tsirelson1980}.
\end{enumerate}
\end{theorem}
\status{published}
\source{P9:thm:chsh_f}

\begin{theorem}[CHSH with density feedback]
\label{P9:thm:chsh_fb}
Let $\CHSH_\mathrm{raw}(f_0,\beta)$ be the raw CHSH of the
$(f_0,\beta)$-model~\eqref{P9:eq:p_fb} in 2D planar geometry at the
standard optimal angles, and let
$\CHSH_\mathrm{rescaled}(f_0,\beta) = \CHSH_\mathrm{raw}(f_0,\beta)/(3/8)$.
 
\begin{enumerate}[label=\emph{(\roman*)}]
\item \emph{(Exact formula.)}
  \begin{align}
    \CHSH_\mathrm{raw}(f_0,\beta)
    &= C_{f_0,\beta}(\pi/8) + C_{f_0,\beta}(3\pi/8),
    \label{P9:eq:chsh_fb_formula}\\
    C_{f_0,\beta}(\psi)
    &= \frac{1}{2\pi}\int_0^{2\pi}
       x_{f_0,\beta}(\varphi)\,x_{f_0,\beta}(\varphi+\psi)\,d\varphi.
    \nonumber
  \end{align}
 
\item \emph{(Effect of $\beta$ at fixed $f_0=2$.)}\label{item:beta_effect}
  At the condensate's own model ($f_0=2$, i.e.\ $p_{2,0}=\pcond$),
  turning on $\beta^*=0.452$ gives:
  \begin{equation}
    \CHSH_\mathrm{rescaled}(2,\,\bst)
    \;\approx\; 2.341,
    \label{P9:eq:chsh_beta_pcond}
  \end{equation}
  slightly \emph{below} the $\beta=0$ value of $2.363$.
  The reduction arises because $\beta>0$ lowers $f_\mathrm{eff}$
  near $\theta=0$, softening the correlation at small mismatch angles.
 
\item \emph{(Paper~VIII limit.)}
  Paper~VIII uses $f_\mathrm{base}=2$ in its convention, equivalent to
  $f_0=4$ here (Remark~\ref{P9:rem:f_convention}).
  At $(f_0,\beta) = (4,\,\bst)$:
  \begin{equation}
    \CHSH_\mathrm{raw}(4,\,\bst) \;\approx\; 0.975,
    \qquad
    \CHSH_\mathrm{rescaled}(4,\,\bst) \;\approx\; 2.60,
    \label{P9:eq:chsh_paper8}
  \end{equation}
  in agreement with the Monte Carlo result of
  Paper~VIII~\cite{Paper8} ($\approx2.611$).
 
\item \emph{(Monotonicity in $f_0$, non-monotonicity in $\beta$.)}\label{item:monotone_fb}
  For fixed $\beta\in[0,1)$, $\CHSH_\mathrm{raw}(f_0,\beta)$ is
  strictly increasing in $f_0$.
  For fixed $f_0$ and increasing $\beta$, the raw CHSH \emph{decreases}
  at small $\beta$ (feedback softens near-aligned correlations) but the
  rescaled CHSH tracks the raw behaviour since the rescale factor $8/3$
  is $\beta$-independent.
\end{enumerate}
\end{theorem}
\status{published}
\source{P9:thm:chsh_fb}

\begin{theorem}[Condensate--photon interaction Hamiltonian]
\label{P9:thm:H_int}
Let $\rho$ be the condensate scalar field with $k$-essence action
\begin{equation}
  S_\mathrm{cond}
  \;=\; \int\!\sqrt{-g}\;P(X)\,d^4x,
  \qquad
  P(X) = \frac{X}{\vp^6(1+\bst X)},
  \qquad
  X = g^{ab}\nabla_{\!a}\rho\,\nabla_{\!b}\rho.
  \label{P9:eq:S_cond_H}
\end{equation}
Let $A_\mu$ be the electromagnetic gauge field.
The minimal coupling of $A_\mu$ to the condensate Noether current is:
\begin{equation}
  H_\mathrm{int}
  \;=\; e\int J^\mu A_\mu\,d^3x,
  \qquad
  J^\mu \;=\; 2P'(X)\,\nabla^\mu\rho,
  \label{P9:eq:H_int}
\end{equation}
where $P'(X) = dP/dX = [\vp^6(1+\bst X)^2]^{-1}$ is the kinetic
susceptibility of the condensate.

Evaluated at the saddle-point background $(\rho_0, X_0=|\nabla\rho_0|^2)$
with the condensate gradient oriented at in-plane angle $\phi_n$, and
coupled to a linearly polarized photon with polarization vector
$\varepsilon^\mu(\alpha) = (0,\cos\alpha,\sin\alpha,0)$:
\begin{equation}
  \langle\phi_n|H_\mathrm{int}|\alpha\rangle_\mathrm{phot}
  \;=\; g_\mathrm{eff}\,\cos(\phi_n - \alpha),
  \label{P9:eq:H_int_matrix}
\end{equation}
where the effective coupling constant is
\begin{equation}
  g_\mathrm{eff}
  \;=\; \frac{2\,e\,A_0\,|\nabla\rho_0|}{\vp^6\,(1+\bst X_0)^2}.
  \label{P9:eq:g_eff}
\end{equation}
The coupling is suppressed by the same factor $\vp^6$ that appears in the
Bell suppression law $S(\theta)=\sin^4\theta/\vp^6$:
both arise from the single $k$-essence action~\eqref{P9:eq:S_cond_H}.
\end{theorem}
\status{published}
\source{P9:thm:H_int}

\begin{remark}[Shared $\vp^6$ denominator: Bell and Malus from one action]
\label{P9:rem:phi6_connection}
The coupling~\eqref{P9:eq:g_eff} is suppressed by $\vp^6\approx17.9$,
the same factor in the Bell suppression law $S(\theta)=\sin^4\theta/\vp^6$.
Both arise from the single action~\eqref{P9:eq:S_cond_H}: the Noether current
$J^\mu=2P'(X)\nabla^\mu\rho$ that controls the photon coupling is the same
object whose square $X=|\nabla\rho|^2$ drives the suppression.
This is the microscopic unity of the Bell and Malus phenomena in the condensate.
\end{remark}
\status{published}
\source{P9:rem:phi6_connection}

\begin{theorem}[Malus's law from the condensate interaction]
\label{P9:thm:malus}
The normalised transition probability for condensate orientation $|\phi_n\rangle$
to produce a photon detected at polarizer angle $\alpha$ is:
\begin{equation}
  T(\phi_n,\alpha)
  \;=\; \frac{|\langle\phi_n|H_\mathrm{int}|\alpha\rangle|^2}
             {|\langle\phi_n|H_\mathrm{int}|\alpha\rangle|^2
              +|\langle\phi_n|H_\mathrm{int}|\alpha{+}\tfrac{\pi}{2}\rangle|^2}
  \;=\; \cos^2(\phi_n-\alpha).
  \label{P9:eq:malus_derived}
\end{equation}
The signed Stokes observable eigenvalue is
$A^\mathrm{Stokes}(\alpha)|\phi_n\rangle = (2T-1)|\phi_n\rangle
= \cos(2(\phi_n-\alpha))|\phi_n\rangle$,
the $m=2$ Fourier mode of $S^1$.
\end{theorem}
\status{published}
\source{P9:thm:malus}

\begin{theorem}[Photon Stokes observable recovers Tsirelson exactly]
\label{P9:thm:stokes}
For a condensate photon singlet with the Stokes observable~\eqref{P9:eq:stokes_op} and uniform ensemble over $\phi_n$:
\begin{equation}
  E^\mathrm{Stokes}(\alpha,\beta)
  \;=\; -\frac{1}{2\pi}\int_0^{2\pi}
  \cos\bigl(2(\varphi-\alpha)\bigr)\,\cos\bigl(2(\varphi-\beta)\bigr)\,d\varphi
  \;=\; -\frac{1}{2}\cos\bigl(2(\alpha-\beta)\bigr).
  \label{P9:eq:E_stokes}
\end{equation}
At the optimal measurement angles $(a,a',b,b')=(0,\pi/4,\pi/8,-\pi/8)$,
after the standard Stokes normalisation $(\times2)$ (where $E_\mathrm{QM}=-\cos(2\psi)$
is normalised to amplitude 1, hence $E^\mathrm{Stokes} = \tfrac{1}{2}E_\mathrm{QM}$),
the physical CHSH value is:
\begin{equation}
  \boxed{\CHSH^\mathrm{Stokes}
  = 2\times\CHSH^\mathrm{Stokes}_\mathrm{raw}
  = 2\sqrt{2}
  \equiv \Tsirelson},
  \label{P9:eq:stokes_tsirelson}
\end{equation}
\end{theorem}
\status{published}
\source{P9:thm:stokes}

\begin{remark}[The Parseval factor is the averaging effect]
\label{P9:rem:parseval_averaging}
The factor $\tfrac{1}{2}$ in~\eqref{P9:eq:E_stokes} is the Parseval factor
for the $m=2$ mode of $L^2(S^1)$.
A photon polarizer selects only the $m=2$ Fourier component of the condensate
orientation distribution; the ensemble average over $\phi_n$ produces this
Parseval factor automatically.
The QM normalisation $\times2$ undoes it, recovering the full amplitude.
Photon Bell experiments \emph{do} perform an ensemble average over $\phi_n$,
and the condensate model provides the geometric explanation of this factor.
\end{remark}
\status{published}
\source{P9:rem:parseval_averaging}

\begin{remark}[The two CHSH ceilings]
\label{P9:rem:two_ceilings}
The condensate model has \emph{two distinct CHSH ceilings}, depending
on the observable:
\begin{center}
\begin{tabular}{llccc}
\toprule
Particle & Observable & Parseval & Rescaling & Ceiling \\
\midrule
Condensate lepton & $\xcond(\theta)$, pcond, $f{\to}\infty$ & $3/8$ & $\times8/3$ & $8/3\approx2.667$ \\
Photon (Stokes)    & $\cos(2\theta)$, Malus, $m{=}2$ (Thm~\ref{P9:thm:malus})        & $1/2$ & $\times2$   & $2\sqrt{2}\approx2.828$ \\
\bottomrule
\end{tabular}
\end{center}
Both ceilings are $\text{CHSH}_\mathrm{raw}\times(1/\text{Parseval factor})$.
Photon Bell experiments~\cite{Hensen2015,Giustina2015,Shalm2015,Poh2015} observe
CHSH $\to2\sqrt{2}$: correct condensate prediction for spin-1/Malus.
The pcond/$8/3$ ceiling is the condensate lepton sector prediction.
The structural Tsirelson gap $2\sqrt{2}-8/3\approx0.161$ is the difference
between the photon and lepton ceilings, a prediction for lepton Bell tests.
\end{remark}
\status{published}
\source{P9:rem:two_ceilings}

\begin{theorem}[Photon winding number from the $\mathbb{RP}^1$ double cover]
\label{P9:thm:m_phot}
The photon polarization observable in the condensate is the $m=2$
Fourier mode of $S^1$, uniquely determined by two constraints:
\begin{enumerate}[label=(\roman*),noitemsep]
\item \emph{(Malus's law, Theorem~\ref{P9:thm:malus})} the transmission
  probability at mismatch angle $\theta=\phi_n-\alpha$ is
  $T(\theta)=\cos^2\theta$, derived from the condensate Noether current;
\item \emph{($\pi$-periodicity)} the observable is invariant under
  $\alpha\mapsto\alpha+\pi$ (rotating the polarizer by $180°$ restores the same measurement).
\end{enumerate}
The eigenvalue of the Stokes observable on $|\phi_n\rangle$ is
$\cos(2(\phi_n-\alpha))$, which is the unique function of $\phi_n-\alpha$
satisfying (i) and (ii).
\end{theorem}
\status{published}
\source{P9:thm:m_phot}

\begin{remark}[Angle doubling: polarizer angle to condensate angle]
\label{P9:rem:winding}
A photon polarizer rotated by $\psi$ corresponds to a condensate
orientation angle of $2\psi$, because photon polarization lives on
$\mathbb{RP}^1 = S^1/\mathbb{Z}_2$ (period $\pi$) while the condensate
lives on $S^1$ (period $2\pi$): the double cover gives winding number
$m_\mathrm{phot}=2$ (Theorem~\ref{P9:thm:m_phot}).
So the effective correlation for a photon pair is
$E_\mathrm{phot}(\psi) = -x_f(2\psi)$ in the projective limit,
or $E_\mathrm{phot}(\psi) = -C_f^{(m=2)}(\psi)$ in general.
This is the physical content of Conjecture~\ref{P9:conj:angle_doubling}.
\end{remark}
\status{published}
\source{P9:rem:winding}

\begin{proposition}[Projective-measurement limit: pcond lepton ceiling]
\label{P9:prop:projective}
For the pcond observable $x_f$ (the condensate's own spin-$\tfrac{1}{2}$
$\HH$-valued model), as $f\to\infty$ the correlation function approaches
a triangular autocorrelation:
\begin{equation}
  C_\infty(\psi) \;\equiv\; \lim_{f\to\infty} C_f(\psi)
  \;=\; \frac{1}{2\pi}\int_0^{2\pi}
  \mathrm{sgn}(\cos\varphi)\,\mathrm{sgn}(\cos(\varphi+\psi))\,d\varphi
  \;=\; 1 - \frac{2|\psi|}{\pi},
  \quad |\psi|\leq\pi.
\end{equation}
At the standard optimal angles:
$C_\infty(\pi/8)=3/4$, $C_\infty(3\pi/8)=1/4$, giving
\begin{equation}
  \CHSH_\mathrm{raw}(\infty) = 1,
  \qquad
  \CHSH_\mathrm{rescaled}(\infty) = \frac{8}{3} \approx 2.667.
  \label{P9:eq:chsh_projective}
\end{equation}
The residual gap from Tsirelson's bound is
\begin{equation}
  \boxed{\Tsirelson - \frac{8}{3}
  \;=\; 2\sqrt{2} - \frac{8}{3}
  \;\approx\; 0.161},
  \label{P9:eq:structural_gap}
\end{equation}
This is the \emph{pcond Tsirelson gap}: the residual within the
$x_f$ observable from the $\sin^4/\vp^6$ suppression geometry.
It does \emph{not} apply to photon polarization experiments,
which use the Stokes observable
(Theorem~\ref{P9:thm:stokes} and Remark~\ref{P9:rem:two_ceilings}).
\end{proposition}
\status{published}
\source{P9:prop:projective}

\begin{remark}[Scope of Proposition~\ref{P9:prop:projective}]
\label{P9:rem:projective_scope}
The ceiling $\CHSH_\mathrm{rescaled}\leq8/3$ in Proposition~\ref{P9:prop:projective}
applies only to the pcond/$x_f$ observable with $\times8/3$ rescaling.
Photon experiments use the Stokes observable ($\times2$, ceiling $2\sqrt{2}$).
\end{remark}
\status{published}
\source{P9:rem:projective_scope}

\begin{proposition}[Incoherent average and the pcond ceiling]
\label{P9:prop:incoherence}
Let $\rho_\mathrm{classic} = (1/2\pi)\int|\phi_n\rangle\langle\phi_n|_A\otimes|\phi_n{+}\pi\rangle\langle\phi_n{+}\pi|_B\,d\phi_n$ be the mixed state arising from incoherent averaging over $\hat{n}$, and with pcond measurement $A(\alpha)|\phi_n\rangle = x_f(\phi_n-\alpha)|\phi_n\rangle$.
The correlation it predicts is:
\begin{equation}
  E_\mathrm{class}(\alpha,\beta)
  \;=\; -C_f(\alpha-\beta)
  \;=\; -\frac{1}{2\pi}\int_0^{2\pi}
  x_f(\varphi-\alpha)\,x_f(\varphi-\beta)\,d\varphi,
  \label{P9:eq:E_class}
\end{equation}
with ceiling $\CHSH_\mathrm{raw}\to1$ and $\CHSH_\mathrm{rescaled}\to8/3\approx2.667$
as $f\to\infty$ (Proposition~\ref{P9:prop:projective}).
This is strictly below Tsirelson's bound.
\end{proposition}
\status{published}
\source{P9:prop:incoherence}

\begin{proposition}[Phase invariance: winding phases cannot break the pcond ceiling]
\label{P9:prop:phase_invariance}
For any winding number $m_0\in\ZZ$, the state~\eqref{P9:eq:qsinglet}
gives the \emph{same} correlation~\eqref{P9:eq:E_class} as the classical
mixed state $\rho_\mathrm{classic}$, provided the measurement operators
are diagonal in the $|\phi\rangle$ basis:
$A(\alpha)|\phi\rangle = x_f(\phi-\alpha)|\phi\rangle$.For pcond measurement operators $A(\alpha)|\phi\rangle = x_f(\phi-\alpha)|\phi\rangle$
and $B(\beta)|\phi+\pi\rangle = x_f((\phi+\pi)-\beta)|\phi+\pi\rangle$,
the correlation from $|\Psi_\mathrm{cond}\rangle$ is:
 \begin{equation}
  E(\alpha,\beta)
  = \langle\Psi_\mathrm{cond}|A(\alpha)\otimes B(\beta)|\Psi_\mathrm{cond}\rangle
  = -C_f(\alpha-\beta),
\end{equation}
independent of $m_0$.
\end{proposition}
\status{published}
\source{P9:prop:phase_invariance}

\begin{lemma}[Anti-symmetry of $x_f$]
\label{P9:lem:antisymmetry}
For all $\theta\in(0,\pi)$ and $f>0$:
\begin{equation}
  x_f(\pi-\theta) = -x_f(\theta).
  \label{P9:eq:antisymmetry}
\end{equation}
\end{lemma}
\status{published}
\source{P9:lem:antisymmetry}

\begin{definition}[Signed correlation functions]
For a spin-$\tfrac{1}{2}$ measurement at mismatch angle $\theta$,
define the \emph{signed} outcome ($\pm1$ valued) correlation functions:
\begin{align}
  \xcond(\theta) &= 2\pcond(\theta) - 1
  = \boxed{\frac{2\cos\theta}{1+\cos^2\theta}}, \label{P9:eq:xcond}\\
  x_\mathrm{QM}(\theta) &= 2\pQM(\theta) - 1 = \cos\theta. \label{P9:eq:xQM}
\end{align}
\end{definition}
\status{published}
\source{P9:eq:xcond}

\begin{remark}[Stokes vs pcond observables]
\label{P9:rem:stokes_vs_pcond}
Both observables map $|\phi_n\rangle$ to itself (they are diagonal)
but with different eigenvalue functions:
\begin{itemize}[noitemsep]
\item $\xcond(\theta) = 2\cos\theta/(1+\cos^2\theta)$: the condensate
  pcond observable ($f=2$, $\HH$-valued). Runs from $+1$ (aligned)
  through 0 (perpendicular) to $-1$ (anti-aligned) in a non-cosine shape.
\item $\cos(2\theta) = 2\cos^2\theta-1$: the Stokes observable (Malus,
  spin-1). Runs from $+1$ at $\theta=0$ through 0 at $\theta=\pi/4$
  to $-1$ at $\theta=\pi/2$, then back to $+1$ at $\theta=\pi$.
\end{itemize}
At $\theta=0$ both give $+1$.
At $\theta=\pi/4$: $\xcond(\pi/4)\approx0.943$ while $\cos(\pi/2)=0$ ---
the pcond observable is nearly maximal where Stokes vanishes.
This functional difference is what gives the different CHSH ceilings.
The Stokes eigenvalue $\cos(2\theta)$ selects exactly the $m=2$ Fourier
coefficient of the condensate distribution; no other harmonics contribute
to the Parseval autocorrelation.
\end{remark}
\status{published}
\source{P9:rem:stokes_vs_pcond}

\begin{definition}[$\HH$-valued condensate amplitude]
\label{P9:def:H_amplitude}
The quaternionic condensate amplitude at mismatch angle $\theta$ is defined by
\begin{equation}
  |\psi_\HH(\theta)|_\HH^2
  \;\equiv\;
  \frac{e^{-\sin^4\theta/\vp^6}}{Z(\theta)},
  \label{P9:eq:psi_H}
\end{equation}
where $Z(\theta) = e^{-\sin^4\theta/\vp^6}+e^{-\cos^4\theta/\vp^6}$
is the two-outcome partition function, so that $|\psi_\HH(+,\theta)|_\HH^2
+ |\psi_\HH(-,\theta)|_\HH^2 = 1$.
\end{definition}
\status{published}
\source{P9:def:H_amplitude}

\begin{definition}[$f$-$\beta$ condensate probability]
\label{P9:def:p_f_beta}
For measurement-strength parameter $f_0\geq1$ and density-feedback
coupling $\beta\geq0$, define the direction-dependent effective exponent
\begin{equation}
  f_\mathrm{eff}(\theta;\,f_0,\beta)
  \;\equiv\; \frac{f_0}{1 + \beta\,\cos\theta},
  \label{P9:eq:f_eff}
\end{equation}
and the corresponding probability and signed correlation:
\begin{equation}
  p_{f_0,\beta}(\theta)
  \;\equiv\; \frac{1}{1+\tan^{2f_\mathrm{eff}(\theta)}(\theta/2)},
  \qquad
  x_{f_0,\beta}(\theta) \;=\; 2p_{f_0,\beta}(\theta)-1.
  \label{P9:eq:p_fb}
\end{equation}
At $\beta=0$: $p_{f_0,0} = p_{f_0}$ (eq.~\eqref{P9:eq:p_f}).
At $f_0=2$: $p_{2,0} = \pcond$ (the condensate's own model).
\end{definition}
\status{published}
\source{P9:def:p_f_beta}

\begin{definition}[Quantum condensate singlet]
\label{P9:def:quantum_singlet}
The quantum-coherent pair state for the condensate is the angular singlet:
\begin{equation}
  |\Psi_\mathrm{cond}\rangle
  = \frac{1}{\sqrt{2\pi}}
  \int_0^{2\pi} e^{im_0\phi}\,|\phi\rangle_A\,|\phi+\pi\rangle_B\,d\phi,
  \label{P9:eq:qsinglet}
\end{equation}
where $m_0\in\ZZ$ is the \emph{winding number} of the pair, set by the Hopf
charge $Q=10$ and the particle species.
The state assigns amplitude $e^{im_0\phi}/\sqrt{2\pi}$ to the condensate
being oriented at angle $\phi$ for particle~A and anti-aligned ($\phi+\pi$) for
particle~B, summed coherently over all orientations.
\end{definition}
\status{published}
\source{P9:def:quantum_singlet}

\begin{definition}[Winding-coherent measurement operator]
\label{P9:def:winding_measurement}
The winding-coherent measurement operator for polarization angle $\alpha$
at winding number $m$ is:
\begin{equation}
  \hat{A}_m(\alpha)
  \;\equiv\; 2|\alpha\rangle_m{}_m\langle\alpha| - \hat{I},
  \label{P9:eq:A_winding}
\end{equation}
with matrix element
$\langle\phi|\hat{A}_m(\alpha)|\phi'\rangle
= (1/\pi)\cos(m(\phi-\phi'))\,\hat{x}_f(\phi-\alpha)$
where $\hat{x}_f$ is the effective correlation weight.
\end{definition}
\status{published}
\source{P9:def:winding_measurement}

\begin{definition}[Photon Stokes observable]
\label{P9:def:stokes}
A photon polarizer at angle $\alpha$ transmits a photon with probability
$\cos^2(\phi_n-\alpha)$ and reflects it with probability $\sin^2(\phi_n-\alpha)$ coupling to the condensate via Malus's law, where $\phi_n$ is the condensate orientation at the source. The signed observable, \emph{Stokes observable}, (eigenvalue $+1$ for transmitted, $-1$ for reflected) for photon polarization at angle $\alpha$ is:

\begin{equation}
  A^\mathrm{Stokes}(\alpha) \;:\; |\phi_n\rangle
  \;\longmapsto\; \bigl(2\cos^2(\phi_n-\alpha)-1\bigr)|\phi_n\rangle
  \;=\; \cos\bigl(2(\phi_n-\alpha)\bigr)|\phi_n\rangle,
  \label{P9:eq:stokes_op}
\end{equation}

i.e.\ eigenvalue $\cos(2(\phi_n-\alpha))\in[-1,+1]$ on each condensate
orientation $|\phi_n\rangle$.
The identity $2\cos^2\theta-1=\cos(2\theta)$ shows this is the
\emph{pure $m=2$ Fourier mode} of the condensate circle.
That $T(\phi_n,\alpha)=\cos^2(\phi_n-\alpha)$ is the physical transmission
probability is proved from the condensate Noether current in
Theorem~\ref{P9:thm:malus}; that $m=2$ is the unique winding number
satisfying Malus's law and $\pi$-periodicity is proved in
Theorem~\ref{P9:thm:m_phot}.

\end{definition}
\status{published}
\source{P9:def:stokes}

\begin{conjecture}[Quantisation of the Hopfion condensate --- \textcolor{darkgreen}{proved in Paper~X~\cite{Paper10}}]
\label{P9:conj:quantisation}
The canonical quantisation of the Hopfion condensate action
$E_\mathrm{fb}[f;\beta]$~\cite{Paper6,Paper7} produces a
Hilbert space $\mathcal{H}=L^2(\mathcal{C}_Q^\mathrm{red},d\mu)$ where
$\mathcal{C}_Q^\mathrm{red}$ is the Derrick-reduced space of topological
field configurations with Hopf charge $Q=10$, and $d\mu$ is the Liouville
measure of the condensate symplectic structure.
The Born rule $P=|\psi|^2$ follows from the Liouville measure being
the unique $\mathrm{U}(1)$-invariant measure on $\mathcal{C}_Q^\mathrm{red}$,
with $\mathrm{U}(1)$ the symmetry group of the phase $\Phi = Q\cdot\omega$.
\end{conjecture}
\status{open}
\source{P9:conj:quantisation}

\begin{remark}[Resolution in Paper~X]
\label{P9:rem:conj_resolved}
Conjecture~\ref{P9:conj:quantisation} is fully resolved in
Paper~X~\cite{Paper10} (Theorem~7.1 of that paper)
The key additional ingredient supplied by Paper~X is the identification
of the Derrick scale instability as the kernel of the symplectic form
$\Omega$, whose quotient yields the non-degenerate reduced space
$\mathcal{C}_Q^\mathrm{red}$.
Paper~X further proves the golden-ratio energy spectrum
$E_n = \vp^{2n}E_0$, the Macfarlane--Biedenharn $q$-oscillator algebra
with $q=\vp^2$ (with $q$-numbers $[n]_{\vp^2}=F_{2n}$, the even Fibonacci
numbers), and the complete Bekenstein--Hawking thermodynamics
$S_\mathrm{Wald}=A_H/(4G_N)$ (tree level) and $\delta S_\mathrm{1-loop}=0$
(one loop).
\end{remark}
\status{published}
\source{P9:rem:conj_resolved}

\begin{conjecture}[Post-quantum regime via pcond coupling]
\label{P9:conj:angle_doubling}
If a future experiment couples to the condensate via the \emph{pcond}
observable (not Malus's law) for a spin-1 particle with $\pi$-periodic
polarization, the effective mismatch angle doubles. For a photon pair (spin-1) in the condensate, the physical correlation function is:
\begin{equation}
  E_\mathrm{phot}(\alpha,\beta)
  = -\xcond(2(\alpha-\beta))
  = -\frac{2\cos(2\psi)}{1+\cos^2(2\psi)},
  \quad \psi = \alpha-\beta.
  \label{P9:eq:E_phot_conjecture}
\end{equation}

\begin{remark}
\label{P9:rem:poh}
Standard quantum mechanics predicts $E_\mathrm{QM}(\psi) = -\cos(2\psi)$
for a photon singlet; the conjectured condensate
formula~\eqref{P9:eq:E_phot_conjecture}
gives $\xcond(2\psi)\geq\cos(2\psi)$ (proved below), so the condensate
correlation is \emph{stronger} than QM at every angle.
This is distinct from the standard photon Bell test, where the Stokes/Malus
observable is used and the condensate correctly predicts $\CHSH=\Tsirelson$
(Theorem~\ref{P9:thm:stokes}): the Poh et al.~\cite{Poh2015} result
$S=2.82759\pm0.00051$ is exactly the Stokes/Malus prediction and is
\emph{not} in tension with the condensate model.
The post-quantum regime~\eqref{P9:eq:E_phot_conjecture} requires a
\emph{different} coupling to the condensate field --- one in which the
measurement back-action is proportional to the pcond (quaternionic)
weight rather than the Malus transmission probability proved in
Theorem~\ref{P9:thm:H_int} and Theorem~\ref{P9:thm:malus}.
\end{remark}

Since
\begin{equation}
  \xcond(2\psi) \;\geq\; \cos(2\psi)
  \quad\text{for all }\psi\in[0,\pi/2],
  \label{P9:eq:xcond_geq_cos}
\end{equation}
(proved below) and $4\xcond(\pi/4)=8\sqrt{2}/3$,
the CHSH under the pcond coupling at optimal angles
$(a,a',b,b')=(0,\pi/4,\pi/8,-\pi/8)$ would be:
\begin{equation}
  \boxed{\CHSH_\mathrm{cond,phot}
  \;=\; 4\,\xcond(\pi/4)
  \;=\; \frac{8\sqrt{2}}{3}
  \;\approx\; 3.771},
  \label{P9:eq:CHSH_postquantum}
\end{equation}
which exceeds Tsirelson's bound $2\sqrt{2}\approx2.828$.
This is a post-quantum prediction contingent on realizing the pcond coupling
for a spin-1 particle: a physically distinct scenario from the standard
Stokes/Malus photon Bell test.
\end{conjecture}
\status{open}
\source{P9:conj:angle_doubling}

\begin{remark}[No-signaling and the post-quantum window]
\label{P9:rem:nosignaling}
A correlation $E(\alpha,\beta)$ that exceeds the Tsirelson bound is
called \emph{super-quantum} or \emph{post-quantum} (Popescu--Rohrlich box).
Such correlations do not violate no-signaling provided
$\int E(\alpha,\beta)\,d\alpha = 0$ for all $\beta$ (marginal independence).
The conjectured $E_\mathrm{phot} = -\xcond(2\psi)$ satisfies this:
$\int_0^{2\pi}\xcond(2(\phi-\beta))\,d\phi = 0$ by the odd symmetry of
$\xcond$ over a full period. It remains to be seen whether it can be realized without violating local causality.
\end{remark}
\status{open}
\source{P9:rem:nosignaling}

\begin{align}
      \CHSH_\mathrm{raw}(f)
      &= \bigl|C_f(\pi/8) + C_f(3\pi/8)\bigr|,
      \\
      C_f(\psi)
      &= \frac{1}{2\pi}\int_0^{2\pi} x_f(\varphi)\,x_f(\varphi+\psi)\,d\varphi,
      \nonumber
    \end{align}

\begin{equation}
  \CHSH_\mathrm{raw}
  = \bigl|E_{ab}+E_{ab'}\bigr|
  = C(\pi/8)+C(3\pi/8)
  \approx 0.886.
  \label{P9:eq:CHSH_raw_value}
\end{equation}

\begin{equation}
  \CHSH = 2\bigl[|E(\pi/8)+E(3\pi/8)| + |E(\pi/8)-E(\pi/8)|\bigr]\big|_\mathrm{opt}
  = 4|E(\pi/8)|.
  \label{P9:eq:CHSH_simplified}
\end{equation}

\begin{equation}
  C(\psi)
  \;=\; \sum_{k=0}^{\infty}\frac{A_{2k+1}^2}{2}\cos((2k+1)\psi)
  \;=\; \frac{(4-2\sqrt{2})^2}{2}\,
        \mathrm{Re}\!\left[\frac{e^{i\psi}}{1-(3-2\sqrt{2})^2 e^{2i\psi}}\right],
  \label{P9:eq:C_exact}
\end{equation}

\begin{remark}[Winding autocorrelation: general $(f,m)$ framework]
\label{P9:rem:quantum_coherent}
For a condensate singlet~\eqref{P9:eq:qsinglet} measured at winding number $m$,
the correlation function generalises to:
\begin{equation}
  E_\mathrm{QC}(\alpha,\beta;f,m)
  \;=\; -C_f^{(m)}(\alpha-\beta),
  \label{P9:eq:E_QC}
\end{equation}
where the \emph{winding autocorrelation} is
\begin{equation}
  C_f^{(m)}(\psi)
  \;\equiv\; \frac{1}{2\pi}\int_0^{2\pi}
  x_f^{(m)}(\varphi)\,x_f^{(m)}(\varphi+\psi)\,d\varphi,
  \qquad
  x_f^{(m)}(\theta) \;\equiv\; x_f(m\theta/2).
  \label{P9:eq:C_winding}
\end{equation}
At $m=1$: $x_f^{(1)}=x_f$ (condensate baseline).
At $m=2$ (photon): $x_f^{(2)}(\theta)=x_f(\theta)$; the angle-doubling enters
through the PHYSICAL angle $\psi=\alpha-\beta$ being mapped to condensate angle $2\psi$,
as derived from the $\mathbb{RP}^1\xrightarrow{2:1}S^1$ double cover
(Remark~\ref{P9:rem:winding} and Theorem~\ref{P9:thm:m_phot}).
At $m=2$ (photon, Stokes/Malus coupling): this reduces exactly to
$E^\mathrm{Stokes}(\psi)=-\tfrac{1}{2}\cos(2\psi)$ of Theorem~\ref{P9:thm:stokes}.
\end{remark}
\status{published}
\source{P9:rem:quantum_coherent}

\begin{equation}
  E_\mathrm{raw}(\psi)
  \;=\; -\frac{1}{2\pi}\int_0^{2\pi} \xcond(\varphi)\xcond(\varphi+\psi)\,d\varphi
  \;=\; -C(\psi),
  \label{P9:eq:E_raw}
\end{equation}

\begin{equation}
  J^\mu
  = \frac{\partial(\sqrt{-g}\,P)}{\partial(\nabla_\mu\rho)}
  = 2P'(X)\,\nabla^\mu\rho.
  \label{P9:eq:J_Noether}
\end{equation}

\begin{equation}
  J^\mu\varepsilon_\mu(\alpha)
  = \frac{2|\nabla\rho_0|}{\vp^6(1+\bst X_0)^2}
  \cos(\phi_n-\alpha).
  \label{P9:eq:J_eps}
\end{equation}

\begin{equation}
  J^\mu\big|_{\rho_0}
  = \frac{2|\nabla\rho_0|}{\vp^6(1+\bst X_0)^2}\,
  (\cos\phi_n,\;\sin\phi_n,\;0,\;0).
  \label{P9:eq:J_saddle}
\end{equation}

\begin{equation}
  P_\HH(\theta)
  = \frac{|\langle +|v(\theta)\rangle|_\HH^4}
         {|\langle +|v(\theta)\rangle|_\HH^4 + |\langle -|v(\theta)\rangle|_\HH^4}
  = \frac{\cos^4(\theta/2)}{\cos^4(\theta/2)+\sin^4(\theta/2)}
  = \pcond(\theta).
  \label{P9:eq:PH_is_Pcond}
\end{equation}

\begin{equation}
  P'(X)
  = \frac{1}{\vp^6(1+\bst X)^2}.
  \label{P9:eq:Pprime}
\end{equation}

\begin{equation}
  P(a) = \mathrm{Tr}(\rho\, \Pi_a)
  \label{P9:eq:born_standard}
\end{equation}

\begin{align}
    \CHSH_\mathrm{raw}(f_0,\beta)
    &= C_{f_0,\beta}(\pi/8) + C_{f_0,\beta}(3\pi/8),
    \\
    C_{f_0,\beta}(\psi)
    &= \frac{1}{2\pi}\int_0^{2\pi}
       x_{f_0,\beta}(\varphi)\,x_{f_0,\beta}(\varphi+\psi)\,d\varphi.
    \nonumber
  \end{align}

\begin{equation}
  \tan\!\left(\frac{\vartheta}{2}\right)
  \;=\; \tan^2\!\left(\frac{\theta}{2}\right).
  \label{P9:eq:hopf_angle}
\end{equation}

\begin{equation}
  \eta(q_0+q_1\mathbf{i}+q_2\mathbf{j}+q_3\mathbf{k})
  \;=\; \frac{q_1+q_2\mathbf{i}}{q_0+q_3\mathbf{i}} \;\in\; \CC\cup\{\infty\}.
  \label{P9:eq:hopf_map}
\end{equation}

\begin{equation}
  \pQM(\theta) \;=\; \cos^2(\theta/2).
  \label{P9:eq:pQM_def}
\end{equation}

\begin{equation}
  p_f(\theta)
  \;=\; \frac{\cos^{2f}(\theta/2)}{\cos^{2f}(\theta/2)+\sin^{2f}(\theta/2)}
  \;=\; \frac{1}{1+\tan^{2f}(\theta/2)},
  \label{P9:eq:p_f}
\end{equation}

\begin{remark}[Parametrisation conventions]
\label{P9:rem:f_convention}
The exponent in~\eqref{P9:eq:p_f} is $2f$, so:
\begin{align*}
  f=1:&\quad p_f = \tfrac{1}{1+\tan^2(\theta/2)} = \cos^2(\theta/2)
         \;\text{(standard QM Born rule, $\CC$-valued);}\\
  f=2:&\quad p_f = \tfrac{1}{1+\tan^4(\theta/2)} = \pcond
         \;\text{(quaternionic, $\HH$-valued; the condensate's own model).}
\end{align*}
The prior companion paper~\cite{Paper8} writes its weights as
$w_\pm = [\cos^4(\theta/2)]^{f_\mathrm{base}}$ and $[\sin^4(\theta/2)]^{f_\mathrm{base}}$,
giving $p = 1/(1+\tan^{4f_\mathrm{base}}(\theta/2))$.
Its default $f_\mathrm{base}=2$ therefore corresponds to $f=4$ in
the convention of~\eqref{P9:eq:p_f} (exponent $\tan^8$), \emph{not} $f=2$.
The condensate's own model ($\pcond$, $f=2$ here) corresponds to
$f_\mathrm{base}=1$ in Paper~VIII's convention --- a distinction that
explains the numerical difference between the two papers' CHSH values.
\end{remark}
\status{published}
\source{P9:rem:f_convention}

\begin{align}
  \frac{1}{2\pi}\int_0^{2\pi}\cos(2(\varphi-\alpha))\cos(2(\varphi-\beta))\,d\varphi
  &= \frac{1}{2}\bigl(\cos(2\alpha)\cos(2\beta)+\sin(2\alpha)\sin(2\beta)\bigr) \notag\\
  &= \frac{1}{2}\cos\bigl(2(\alpha-\beta)\bigr).
  \label{P9:eq:parseval_m2}
\end{align}

\begin{align}
  \pcond(\theta) &= \frac{1}{1+t^4}
  \;=\; \pQM\bigl|_{t\,\to\, t^2}, \label{P9:eq:pcond_as_pQM}\\
  \pQM(\theta)   &= \frac{1}{1+t^2}. \label{P9:eq:pQM_tan}
\end{align}

\begin{equation}
  \pcond(\theta)
  \;=\; \frac{\cos^4(\theta/2)}{\cos^4(\theta/2)+\sin^4(\theta/2)}.
  \label{P9:eq:pcond_def}
\end{equation}

\begin{equation}
  |\alpha\rangle_\mathrm{phot}
  \;=\; \frac{1}{\sqrt{2\pi}}
  \int_0^{2\pi} e^{i\,m_\mathrm{phot}\,(\phi-\alpha)}\,|\phi\rangle\,d\phi,
  \label{P9:eq:photon_state}
\end{equation}

\begin{equation}
  \langle\sin^{2f}\theta\rangle_\mathrm{2D}
  \;=\; \frac{1}{\pi}\int_0^{\pi}\sin^{2f}\theta\,d\theta
  \;=\; \frac{\Gamma(f+\tfrac{1}{2})}{\sqrt{\pi}\,\Gamma(f+1)}.
  \label{P9:eq:sin2f_avg}
\end{equation}

\begin{equation}
  \langle\sin^4\theta\rangle_\mathrm{2D}
  = \frac{1}{\pi}\int_0^{\pi}\sin^4\theta\,d\theta
  = \frac{3}{8}.
  \label{P9:eq:sin4_2D_exact}
\end{equation}

\begin{align}
  \xcond(\theta) &= 2\pcond(\theta) - 1
  = \boxed{\frac{2\cos\theta}{1+\cos^2\theta}}, \\
  x_\mathrm{QM}(\theta) &= 2\pQM(\theta) - 1 = \cos\theta. 
\end{align}

